\documentclass[11pt]{article}
\usepackage[T1]{fontenc}
\usepackage{lmodern}
\usepackage[utf8]{inputenc}
\usepackage{amsmath,amssymb,amsthm}
\usepackage[a4paper,margin=28mm]{geometry}
\usepackage{booktabs}
\usepackage[hidelinks]{hyperref}
\hypersetup{
  pdftitle={A Degree-Spread Bound Between Independent Domination and Minimum Maximal Matchings: A Conjecture},
  pdfauthor={Hassine Achour},
  pdfsubject={A graph theory conjecture relating independent domination, minimum maximal matchings, and degree spread}
}

\newtheorem{conjecture}{Conjecture}
\newtheorem{proposition}{Proposition}
\newtheorem{remark}{Remark}
\newcommand{\mustar}{\mu^{\star}}

\title{A Degree-Spread Bound Between Independent Domination and Minimum Maximal Matchings: A Conjecture}
\author{Hassine Achour\\\texttt{hassine.achour@gmail.com}\\
\small Independent Researcher\\
\small\url{https://doi.org/10.5281/zenodo.21806398}}
\date{August 5, 2026}

\begin{document}
\maketitle

\begin{abstract}
For a finite simple graph $G$, let $i(G)$ be its independent domination number, let $\mustar(G)$ be the minimum cardinality of a maximal matching, and let $\Delta(G)$ and $\delta(G)$ be its maximum and minimum degrees. We conjecture that every connected nontrivial graph satisfies
\[
 i(G)-\mustar(G)\leq \Delta(G)-\delta(G).
\]
For regular graphs this specializes to the open TxGraffiti conjecture $i(G)\leq\mustar(G)$. The proposed correction by degree spread is best possible: equality holds for every balanced double star. Connectedness and nontriviality are necessary. We report an exhaustive verification for all $273{,}192$ connected nonisomorphic graphs of orders $2$ through $9$, together with supplementary structured and random tests. These computations are finite evidence only and do not constitute a proof.
\end{abstract}

\section{Definitions and conjecture}
All graphs considered are finite, simple, and undirected. A set of vertices is \emph{independent dominating} if it is both independent and dominating. The minimum cardinality of such a set is the \emph{independent domination number} $i(G)$ \cite{goddard2013}.

A matching is \emph{maximal} if no edge can be added while preserving the matching property. We write $\mustar(G)$ for the minimum cardinality of a maximal matching. This parameter is also called the \emph{saturation number} or the \emph{edge domination number}; equivalently, $\mustar(G)=i(L(G))=\gamma(L(G))$, where $L(G)$ is the line graph \cite{tavakoli2023}. The notation $\mustar(G)$ is used here to avoid confusion with the unrelated extremal parameter $\operatorname{sat}(n,H)$.

The \emph{degree spread} of $G$ is $\Delta(G)-\delta(G)$. We propose the following inequality.

\begin{conjecture}[Degree-spread conjecture]\label{conj:main}
Every finite, simple, connected, nontrivial graph $G$ satisfies
\begin{equation}\label{eq:main}
 \boxed{\quad i(G)\leq \mustar(G)+\Delta(G)-\delta(G).\quad}
\end{equation}
Equivalently,
\[
 i(G)-\mustar(G)\leq\Delta(G)-\delta(G).
\]
\end{conjecture}

If $G$ is $r$-regular, then $\Delta(G)-\delta(G)=0$, and Conjecture~\ref{conj:main} becomes
\[
 i(G)\leq\mustar(G).
\]
This regular-graph inequality is Conjecture~3 in the TxGraffiti retrospective of Davila, Brimkov, and Pepper and is recorded there as open since 2020 \cite{davila2025}. Thus Conjecture~\ref{conj:main} is a proposed extension of that open problem from regular graphs to arbitrary connected graphs, with an additive correction measuring nonregularity. It should not be interpreted as an independent proof strategy for the regular case.

The related theorem of Rautenbach and Volkmann states that $i(G)+\nu(G)\leq |V(G)|$, where $\nu(G)$ is the maximum matching number \cite{rautenbach2002}. Since $\mustar(G)$ is a minimum maximal matching parameter rather than $\nu(G)$, that theorem does not imply \eqref{eq:main}.

\section{Sharpness and necessary hypotheses}
Let $D_k$ be the balanced double star obtained from an edge $xy$ by adjoining $k$ private leaves to $x$ and $k$ private leaves to $y$, where $k\geq1$.

\begin{proposition}\label{prop:sharp}
For every $k\geq1$, equality holds in \eqref{eq:main} for $D_k$.
\end{proposition}

\begin{proof}
The central edge $xy$ is a maximal matching, so $\mustar(D_k)=1$. An independent dominating set can contain one center together with all $k$ leaves adjacent to the other center, giving size $k+1$. Conversely, if an independent dominating set contains neither center, it contains all $2k$ leaves; if it contains one center, independence excludes the other center and domination forces all $k$ leaves adjacent to that other center. Hence $i(D_k)=k+1$. Finally,
\[
 \Delta(D_k)=k+1,\qquad \delta(D_k)=1,
\]
and therefore
\[
 i(D_k)-\mustar(D_k)=k=\Delta(D_k)-\delta(D_k).
\]
\end{proof}

Consequently, the coefficient of the degree-spread term cannot be reduced below $1$.

The nontriviality assumption is necessary: for $K_1$, one has $i(K_1)=1$, $\mustar(K_1)=0$, and $\Delta(K_1)-\delta(K_1)=0$. Connectedness is also necessary. For the disjoint union $2P_4$, additivity across components gives
\[
 i(2P_4)=4,\qquad \mustar(2P_4)=2,
\]
while $\Delta(2P_4)-\delta(2P_4)=1$. Thus $4\nleq 3$.

\section{Computational evidence}
The invariant $i(G)$ was computed exactly by branch-and-bound over independent dominating sets. The value $\mustar(G)$ was computed using the identity $\mustar(G)=i(L(G))$, and every reported witness was checked directly. Small instances were independently cross-checked by transparent subset enumeration. Connected nonisomorphic graphs were generated with \textsc{nauty} 2.8.8 \cite{mckay2014}.

The exhaustive campaign covered every connected nonisomorphic graph separately at each order from $2$ through $9$.

\begin{center}
\begin{tabular}{rrr}
\toprule
Order & Connected graphs & Violations \\
\midrule
2 & 1 & 0 \\
3 & 2 & 0 \\
4 & 6 & 0 \\
5 & 21 & 0 \\
6 & 112 & 0 \\
7 & 853 & 0 \\
8 & 11,117 & 0 \\
9 & 261,080 & 0 \\
\midrule
Total & 273,192 & 0 \\
\bottomrule
\end{tabular}
\end{center}

For each graph, the recorded slack was
\[
 \mustar(G)+\Delta(G)-\delta(G)-i(G).
\]
The exhaustive campaign found no negative slack and found $34$ equality graphs. The slack distribution was $34,668,6827,35304,87635,86071,39421,12679,3826,726,1$ for slacks $0,1,\ldots,10$, respectively.

A supplementary campaign tested $452$ additional graphs: $120$ connected Erd\H{o}s--R\'enyi samples, $60$ random trees, $60$ random connected bipartite graphs, and $212$ balanced or unbalanced double stars. It found no violation. These supplementary tests may overlap structurally with graphs in the exhaustive campaign and are therefore not added to the exhaustive count as distinct instances.

No finite computation establishes Conjecture~\ref{conj:main} for arbitrary order. In particular, since its regular specialization is already open, the computational evidence should be understood only as support for investigating the proposed degree-spread extension.

\section{Status and questions}
The immediate question is to prove or disprove Conjecture~\ref{conj:main}. Natural restricted cases include trees, bipartite graphs, subcubic graphs, and graphs with $\Delta-\delta=1$. Equality classification is also open beyond the families and small examples reported here.

Searches through August 5, 2026 used the terms independent domination, edge domination, minimum maximal matching, saturation number, degree spread, almost-regular graphs, and the common notations for these parameters. No previous occurrence of the exact inequality \eqref{eq:main} was identified. This is a limited, dated priority statement, not a guarantee that an equivalent formulation does not appear in unindexed or differently notated literature.

\section*{Data and code availability}
The source, exact verification code, tests, and machine-readable reports are archived with this note at
\href{https://doi.org/10.5281/zenodo.21806398}{Zenodo}.

\begin{thebibliography}{9}
\bibitem{davila2025}
R. Davila, B. Brimkov, and R. Pepper,
\emph{In Reverie Together: Ten Years of Mathematical Discovery with a Machine Collaborator},
arXiv:2507.17780v1, 2025.
\url{https://arxiv.org/abs/2507.17780v1}.

\bibitem{goddard2013}
W. Goddard and M. A. Henning,
\emph{Independent domination in graphs: A survey and recent results},
Discrete Mathematics 313 (2013), 839--854.
\url{https://doi.org/10.1016/j.disc.2012.11.031}.

\bibitem{mckay2014}
B. D. McKay and A. Piperno,
\emph{Practical graph isomorphism, II},
Journal of Symbolic Computation 60 (2014), 94--112.
\url{https://doi.org/10.1016/j.jsc.2013.09.003}.

\bibitem{rautenbach2002}
D. Rautenbach and L. Volkmann,
\emph{Independent domination and matchings in graphs},
Discrete Mathematics 259 (2002), 325--330.
\url{https://doi.org/10.1016/S0012-365X(02)00304-7}.

\bibitem{tavakoli2023}
M. Tavakoli and T. Do\v{s}li\'c,
\emph{Smallest maximal matchings of graphs},
Hacettepe Journal of Mathematics and Statistics 52(2) (2023), 356--366.
\url{https://doi.org/10.15672/hujms.1095437}.
\end{thebibliography}

\end{document}
